\documentclass[12pt,a4paper]{article}
% ============================================================
%  Structural Persistence Series — Shared Preamble
% ============================================================
\usepackage{fontspec}
\setmainfont{Hiragino Mincho ProN}
\setsansfont{Hiragino Sans}
\setmonofont{Menlo}

\usepackage{iftex}
\ifXeTeX
  \XeTeXlinebreaklocale "ja"
  \XeTeXlinebreakskip=0pt plus 1pt minus 0.1pt
\fi

\usepackage[margin=22mm]{geometry}
\usepackage{amsmath,amssymb,amsthm}
\usepackage{xcolor}
\usepackage{graphicx}
\usepackage{hyperref}
\usepackage{booktabs}
\usepackage{fancyhdr}
% enumitem not available in this TeX installation

% --- Colours ------------------------------------------------
\definecolor{PaperAccent}{HTML}{1F4E79}
\definecolor{PaperGray}{HTML}{51606F}

\hypersetup{
  colorlinks=true,
  linkcolor=PaperAccent,
  urlcolor=PaperAccent,
  citecolor=PaperAccent,
  pdflang=ja
}
\urlstyle{same}

% --- Typography ---------------------------------------------
\setlength{\parindent}{1.5em}
\setlength{\parskip}{0pt}
\setlength{\emergencystretch}{3em}
\linespread{1.08}
\setlength{\abovedisplayskip}{8pt plus 2pt minus 4pt}
\setlength{\belowdisplayskip}{8pt plus 2pt minus 4pt}
\setlength{\abovedisplayshortskip}{6pt plus 2pt minus 2pt}
\setlength{\belowdisplayshortskip}{6pt plus 2pt minus 2pt}

% --- Section label names ------------------------------------
\renewcommand{\abstractname}{要旨}

% --- Theorem-like environments ------------------------------
\theoremstyle{plain}
\newtheorem{proposition}{命題}[section]
\newtheorem{lemma}[proposition]{補題}
\newtheorem{corollary}[proposition]{系}

\theoremstyle{definition}
\newtheorem{definition}{定義}[section]
\newtheorem{assumption}{条件}[section]

\theoremstyle{remark}
\newtheorem*{remark}{注}

% --- Running header -----------------------------------------
\pagestyle{fancy}
\fancyhf{}
\renewcommand{\headrulewidth}{0.4pt}
\fancyhead[L]{\small\textcolor{PaperGray}{\nouppercase{\leftmark}}}
\fancyhead[R]{\small\textcolor{PaperGray}{\thepage}}
\fancypagestyle{plain}{%
  \fancyhf{}
  \renewcommand{\headrulewidth}{0pt}
  \fancyfoot[C]{\small\textcolor{PaperGray}{\thepage}}
}

% --- Title block macro --------------------------------------
\newcommand{\PaperTitleBlock}[3]{%
  % #1 = title, #2 = subtitle, #3 = date
  \thispagestyle{plain}%
  \begin{center}
  {\LARGE\bfseries #1\par}
  \vspace{0.4em}
  {\normalsize #2\par}
  \vspace{1.0em}
  {\large Akihito Sunagawa\par}
  \vspace{0.3em}
  {\normalsize #3\par}
  \end{center}
  \vspace{1.6em}
}

% Legacy 2-arg version (backward compat)
\newcommand{\PaperTitleBlockSimple}[2]{%
  \PaperTitleBlock{#1}{}{#2}
}


\begin{document}

\PaperTitleBlock
  {構造持続と推論性能の劣化}
  {— 論理一貫性維持の記述と最初の実験 —}
  {2026年4月12日}

\begin{abstract}
大規模言語モデルでは、対話が長くなるにつれて矛盾や未整理の更新が蓄積し、自己矛盾なく推論を完了することが難しくなることがある。本稿では、この劣化を「論理一貫性を保って最後まで到達できる推論経路の集合が、制約の蓄積によって縮小していく過程」として記述する。ここで扱うのは推論性能全般ではなく、与えられた前提や規則のもとで論理一貫性を維持できるかという、より限定された問題である。

対話実験では、外部代謝パイプラインを備えた条件（ON）が、備えない条件（OFF）に対して論理一貫性維持率を上回る方向が観測された。ここで ON とは、矛盾する更新を外部で検出し、「旧値 -\textgreater{} 新値」の時間ラベルつき対として整理・保持する条件である。OFF とは、同じ矛盾を注入するが、その整理を行わず、古い情報と新しい情報を未整理のまま混在させる条件である。gemma3:27b（\(n=3\)、ON/OFF/NC の三条件、180 ターン）では、規則＋事実の合算で ON 73.3\%、NC 56.7\%、OFF 21.1\% となった（\(ON vs OFF p = 0.0004\), \(d = 8.80\); \(ON vs NC p = 0.031\), \(d = 2.67\)）。qwen3.5:27b（\(n=1\)、ON/OFF、30 ターン、§8.3 の実装前提を満たす条件、ON の代謝 LLM は Claude Sonnet）でも、全体正答率は ON 64.4\%、OFF 44.4\% となり、同じ方向が確認された。加えて、100ターンの長期実験（qwen3.5:9b、代謝あり条件）では、論理一貫性は長期にわたって安定し、単調な崩壊は観測されなかった。
\end{abstract}

\section{はじめに}\label{ux306fux3058ux3081ux306b}

本稿は、構造持続の最小形式を、推論時の論理一貫性の劣化に適用する。継続学習における破滅的忘却は次稿で扱う。

大規模言語モデルの推論能力は向上しているが、長い対話では、文脈中に矛盾や未整理の更新が蓄積し、前提や規則を安定して保つことが難しくなることがある。ここで問題にしているのは、単にコンテキストが長くなることではない。古い規則と新しい規則、異なる指示、未整理の参照元が混在することで、整合的な推論の道筋そのものが狭まっていく点である。

本稿で扱うのは、推論性能全般ではなく、与えられた前提や規則のもとで自己矛盾なく推論を完了できるかという意味での論理一貫性維持である。崩れとは、単に誤答が出ることではなく、そのために必要な推論経路が減少し、その結果として一貫した推論を安定して保てなくなることを指す。

本稿の目的は、この事態を最小形式で記述することである。中心に置くのは、論理一貫性の成否そのものを直接説明することではない。まず、論理一貫性を保てる推論経路の集合が、制約の蓄積によってどのように縮小するかを記述する。そのうえで、観測される論理一貫性維持率との関係を、仮定を明示した弱い命題として位置づける。

\section{観測量の定義}\label{ux89b3ux6e2cux91cfux306eux5b9aux7fa9}

観測量 Y を次のように定義する。

\begin{definition}[観測量 Y]
$Y = 1$
与えられた論理課題に対し、モデルが自己矛盾なく、
与えられた規則や前提に反しない形で推論を完了できた。

$Y = 0$
それができなかった。

\end{definition}

重要なのは、Y を先に独立の観測量として定義することである。有効経路の集合を用いて Y を定義してはならない。さもなければ、後に与える理論構成が循環するからである。

\section{候補推論経路と有効経路}\label{ux5019ux88dcux63a8ux8ad6ux7d4cux8defux3068ux6709ux52b9ux7d4cux8def}

ある論理課題に対して、モデルが取りうる候補推論経路の全体を P とする。最初の段階では、経路の細かさを次のような粗い単位で固定する。

\begin{itemize}
\tightlist
\item
  前提を選ぶ
\item
  規則を適用する
\item
  中間結論を出す
\item
  最終結論を出す
\end{itemize}

したがって、P は離散的な推論段階列の集合として扱う。トークン単位の内部過程までは、ここでは立ち入らない。たとえば、「最近はラーメンを優先する」「辛い料理は避ける」という条件のもとで食事提案を行う課題なら、前提を選ぶ -\textgreater{} 規則を適用する -\textgreater{} 候補を絞る -\textgreater{} 結論を出す、という段階列を一つの粗い経路として数える。

制約が加わる前に、論理一貫性を保って最後まで到達しうる候補経路の集合を \[
A^{(0)}
\]

とする。\(A^{(0)}\) は、前提を落とさず、規則に反せず、中間段階で自己矛盾を起こさず、最終結論まで到達できるような推論段階列の集合である。

ここでも、\(A^{(0)}\) を Y そのもので定義しないことが重要である。\(A^{(0)}\) はあくまで、論理一貫性を達成しうる候補経路の構成的集合として置かれる。本稿で導入する有効経路集合は理論構成物であり、その直接観測を主張するものではない。経験的内容は、事前に定めた代理指標が基準モデルに対して追加予測力を持つかどうかにかかる。

\section{制約と縮小}\label{ux5236ux7d04ux3068ux7e2eux5c0f}

論理一貫性を崩しうる制約として、少なくとも次を考える。

\begin{enumerate}
\def\labelenumi{\arabic{enumi}.}
\item
  構造的な矛盾 前提や規則そのものが互いに衝突する。
\item
  指示の衝突 「こう答えよ」と「こう答えるな」が競合する。
\item
  時間的な不整合 古い規則と新しい規則が整理されずに混在する。
\item
  参照元どうしの衝突 参照元ごとに異なる規則や事実が未整理で共存する。
\end{enumerate}

\section{未整理の上書き}\label{ux672aux6574ux7406ux306eux4e0aux66f8ux304d}

新しい情報が古い情報を条件づけなしに上書きする。

本稿では、これらの制約は有効経路の集合に作用し、その集合を縮小させるものとして扱う。したがって、制約が順に加わると、 \[
A^{(0)} \supseteq A^{(1)} \supseteq \cdots \supseteq A^{(n)}
\]

という包含列が得られると仮定する。

\begin{enumerate}
\def\labelenumi{\arabic{enumi}.}
\setcounter{enumi}{4}
\tightlist
\item
  有効経路縮小の指数表現
\end{enumerate}

制約が順に加わるたびに、有効経路の集合が \[
A^{(0)}, A^{(1)}, \cdots, A^{(n)}
\]

と縮小していくとする。各段階の縮小率を \[
p_{i} = m(A^{(i)}) / m(A^{(i-1)})
\]

とおく。ここで m は最小形式における測度に対応する量であり、最初の段階では候補推論経路を一様に数える素朴な測り方を仮置きする。

すると、 \[
m(A^{(n)}) = m(A^{(0)}) \times p_1 \times p_2 \times \cdots \times p_{n}
\]

と書ける。

ここで、累積損失を \[
L = \sum[-log p_{i}]
\]

と置けば、 \[
m(A^{(n)}) = m(A^{(0)}) e^{-L}
\]

となる。段階損失を縮小率の対数比として測る関数形そのものは、Paper 1 §3 の対数比の一意性定理（公理系 B1--B4 の帰結）によって公理的に特徴づけられている。本稿ではその特徴づけを推論ドメインの \(p_{i}\) に対してそのまま適用する。

ここで言いたいのは、制約が有効経路を順に削るなら、その残り方は累積損失 L に対して指数型で書ける、ということである。

ただし、この段階では観測量 Y との関係はまだ主張しない。したがって、ここでの指数型はあくまで経路集合の縮小に関する形式的結果であり、最小形式（Paper 1 §3 の対数比の一意性定理および §4 命題1）の本設定への適用である。本稿では単調性の検証にとどまるが、この指数型は将来の定量的予測、たとえば制約の種類ごとの減衰率の比較へ接続するための足場として置かれている。

\section{経路縮小と観測量を結ぶ三つの前提}\label{ux7d4cux8defux7e2eux5c0fux3068ux89b3ux6e2cux91cfux3092ux7d50ux3076ux4e09ux3064ux306eux524dux63d0}

本稿では、経路集合の縮小と観測量 Y とを結ぶために、次の三つの前提を置く。

\begin{assumption}[仮定A]
構造を保てたなら、そのための経路が必要である。

すなわち、$Y = 1$が成立するならば、その出力過程は少なくとも一つ、論理一貫性を保てる経路を通っている。

これは概念的な前提である。ここへの批判は、論理一貫性維持の定義そのもの、あるいは有効経路の置き方に向かう。
\end{assumption}

\begin{assumption}[仮定B]
観測に使う指標 z は、有効経路の残り方を単調に反映する。

たとえば、z が大きくなるほど、有効経路の残りは小さくなる、あるいは少なくともその傾向を持つと考える。候補となる指標としては、明示的な矛盾の数、未整理の上書きの数、範囲づけされていない衝突の数、指示の衝突数、参照元どうしの不整合密度などがある。

本稿の中心実験では、z を主として「未整理の上書きの数」として具体化する。つまり、新しい情報が古い情報を条件づけなしに置き換え、両者の関係が整理されないまま残る度合いを、操作的な代理指標として用いる。

これは経験的な前提である。したがって、反例は「この指標では理論にうまく乗らない」という形で現れる。
\end{assumption}

\begin{assumption}[仮定C]
制約は縮小方向に働く。

すなわち、各制約は、有効経路の集合を広げることはなく、縮めるか、そのままに保つ。
\end{assumption}

\section{弱い命題}\label{ux5f31ux3044ux547dux984c}

仮定A のもとでは、観測される論理一貫性維持率 Y は、有効経路の残り方によって制約されると考えられる。

ただし、本稿で直接に検証したいのは、この上界そのものではない。本稿の中心に置くのは、次の命題である。

\begin{proposition}[主たる命題]
制約の蓄積を反映する指標 z が増えるにつれて、論理一貫性維持率 Y は平均傾向として悪化する。
\end{proposition}

この命題は、有効経路集合の指数的縮小から直接に導かれる強い上界ではなく、方向と傾向に関する主張にとどまる。本稿でこれを弱い命題と呼ぶのはそのためである。ここでの目的は、完全な特徴づけを与えることではなく、経路集合の縮小という構造的記述と、観測量 Y の悪化方向とを結ぶ最初の関係を与えることにある。

本稿の実質的主張は、仮定Bの妥当性にかかっている。すなわち、z が有効経路の残り方を反映する操作的な代理指標として機能するかどうかが、この命題の経験的内容を決定する。

\section{補助命題: 外部代謝介入の効果}\label{ux88dcux52a9ux547dux984c-ux5916ux90e8ux4ee3ux8b1dux4ecbux5165ux306eux52b9ux679c}

§7 までで、制約の蓄積が経路集合を縮小させ、Y の悪化方向を与えることを見た。ここで問えるのは、この縮小が \textbf{整理によって抑制できるか} である。すなわち、累積損失 L が蓄積する過程に対して、介入を挟むことで残存可能性を維持できるか、である。

本節ではこの問いを補助命題として定式化する。補助命題の役割は、最小形式が予測する「L を抑える介入は S を保つ」という含意を、LLM ドメインで操作的に置き直すことである。

\subsection{補助命題の適用範囲}\label{ux88dcux52a9ux547dux984cux306eux9069ux7528ux7bc4ux56f2}

補助命題を述べる前に、本稿が想定する系の条件を明示する。

\begin{assumption}[想定する系の条件]
本補助命題は、次の性質を持つ系に適用される。
$(i)$ 系は、長期的な矛盾解消代謝機能を内部に持たない。すなわち、複数の推論呼び出しにまたがって矛盾情報を自律的に抽出・統合・記録する機構を持たない。
$(ii)$ 系は、prompt に提示された情報の範囲内でしか状態を更新できない。inference 呼び出しの境界を越えて信念を持ち越す機構を持たない。
\end{assumption}

これらの条件を満たす系に対しては、長期一貫性を維持するための整理プロセスは外部からしか与えられない。本稿ではこの外部整理プロセスを代謝と呼ぶ。

現在の主流的な具体例: Transformer 系の stateless deployment 本稿の執筆時点で広く用いられている Transformer 系 LLM の対話運用の多くは、 少なくとも主流的な stateless deployment という形態において、 上の条件 (i)(ii) を満たす。その理由を簡潔に述べる。 第一に、現在一般的な対話推論の運用では、各 inference 呼び出しは 独立であり、前回の呼び出しで得た情報は prompt に明示的に含めない 限り持ち越されない。 第二に、attention は情報の重み付けを動的に行うが、標準的な forward pass には、古い情報を削除したり、最新値へコミットしたり する明示的な信念更新手続きが存在しない。 第三に、Transformer 系モデルは context を同時並行に参照するため、 「T=1 の X」と「T=14 の Y」はどちらも同じ注意の場で競合する。 どちらを優先するかは位置埋め込みのヒントに頼るほかなく、明示的な 時系列整理は行われない。 第四に、chain-of-thought のような in-context 推論は一度の inference 呼び出し内で完結し、次の呼び出しには持ち越されない。これは外部代謝を inference 内部に押し込んだ形であり、構造としては一時的な整理にとどまる。 したがって、少なくとも現在の主流的な stateless deployment における Transformer 系 LLM は、条件 (i)(ii) を満たす系の具体的な一例である。

他のアーキテクチャへの一般化 本補助命題は Transformer 固有の主張ではない。将来、Transformer とは 異なる構造（たとえば recurrent state を持つ系、外部メモリを持つ系、 継続学習を組み込んだ系）が登場したとしても、そのアーキテクチャが 条件 (i)(ii) を満たすならば、本補助命題の対象系となる。逆に、条件 \((i)\) または \((ii)\) を破る系、すなわち内部で長期的な矛盾解消代謝を行える 系は、本補助命題の適用範囲外となる。そのような系の挙動は、本補助 命題の反証ではなく、対象系の境界を明らかにするものとして扱われる。 本補助命題が主張するのは、あくまで「条件 (i)(ii) を満たす系に対して、 外部からの整理プロセスが構造持続を維持する」という条件つきの方向性 であり、アーキテクチャ固有の経験則ではない。

\subsection{補助命題}\label{ux88dcux52a9ux547dux984c}

\begin{lemma}[補助命題]
§8.1 の条件 (i)(ii) を満たす系 M と、記憶統合性に依存するタスクを評価するベンチマーク B について、次が成立する。

同一の制約注入条件下で、外部からの整理プロセス（代謝）を適用した系（ON）は、それを適用しない系（OFF）に比べて、B 上の維持率 Y が高い。

この命題は方向性の主張であり、具体的な数値比や、カテゴリ別効果サイズの順序を含まない。
\end{lemma}

この命題の核心は \textbf{方向} にある。最小形式が予測するのは「L を下げる介入は S を保つ」という関係であり、L が単調増加量である限り、代謝による L 削減はそれに対応する S の差分を生む。異なるモデルや異なるベンチマークで効果の \textbf{大きさ} は変動しうるが、\textbf{方向} は構造的要請から予期される。

\subsection{補助命題の検証と反証}\label{ux88dcux52a9ux547dux984cux306eux691cux8a3cux3068ux53cdux8a3c}

検証の水準 補助命題の検証は、方向性の再現に限る。 \((v1)\) 同じ \((M, B)\) について、条件 (i)(ii) を満たした上で ON の維持率 が OFF の維持率を上回ること \((v2)\) 複数の \((M, B)\) にわたって、\((v1)\) の方向が保たれること

本稿の §9 で述べる実験はこの水準での最初の観察である。同一の実装条件が維持される限り、より多くの \((M, B)\) ペアでこの方向が観察されれば、補助命題の適用範囲が広いことの経験的支持が積み上がる。

反証条件 補助命題の反証には、次のいずれかが必要である。 \((r1)\) 条件 (i)(ii) を満たす \((M, B)\) において、実装の前提が満たされて いるにもかかわらず \(ON \le OFF\)が観測される \((r2)\) モデル能力が向上するにつれて ON-OFF 差が単調に縮小し、 条件 (i)(ii) を満たす系全体でゼロに収束する兆候が見える

\((r2)\) が成立する場合、それは「条件 (i)(ii) を満たす系で外部代謝の価値が漸減している」ことを意味し、補助命題の経験的内容が弱まる。

実装レベルの前提 補助命題が観測可能になるには、実装レベルで次の四つが満たされている 必要がある。これらは理論の内容ではなく、測定の前提である。 \((p1)\) モデル M は prompt に提示された矛盾を認識できる（in-context での矛盾認識能力） \((p2)\) 代謝プロセスは蓄積された情報から矛盾を抽出・記録できる \((p3)\) retrieval パイプラインは、代謝が記録した情報を応答生成時の prompt に実際に届ける \((p4)\) ベンチマーク B の判定装置は、応答の正誤を正確に評価できる。 表層一致ではなく、意味照合による判定が可能であること

これらの前提は独立に検証可能であるべきであり、いずれかが欠ければ方向性の観察は理論の棄却ではなく測定の失敗として扱われる。

\subsection{適用範囲外の系の扱い}\label{ux9069ux7528ux7bc4ux56f2ux5916ux306eux7cfbux306eux6271ux3044}

条件 (i)(ii) を破る系、たとえば次のような架空または将来の系は、本補助命題の対象外である。

\begin{itemize}
\tightlist
\item
  内部に永続的な状態を持ち、推論呼び出しを越えて信念を更新する系
\item
  訓練中に長期一貫性維持のための整理を学習しており、追加の外部介入を 必要としない系
\item
  外部メモリを持つが、そのメモリ更新が系自身の推論の一部として 自律的に行われる系
\end{itemize}

これらの系に対して \(ON > OFF\)が観測されない場合、それは補助命題の反証ではなく、系が条件 (i)(ii) を満たさないことの確認である。本補助命題は、そのような系の挙動を記述する別の理論構築を妨げない。

重要なのは、条件 (i)(ii) が独立に判定可能である限り、適用範囲は事前に決まっており、事後的に「この系は範囲外だった」と言い訳する余地はないという点である。適用範囲と対象系の分離は、補助命題が循環論法に陥らないための構造的保証である。

\section{実験の位置づけと設計}\label{ux5b9fux9a13ux306eux4f4dux7f6eux3065ux3051ux3068ux8a2dux8a08}

実験には二つの役割がある。

第一に、指標 z と Y の単調関係を確かめること。 第二に、この設定では z が有効経路の残り方を反映する指標として妥当であることを支持することである。

本稿では、z を主として「未整理の上書きの数」として具体化し、その最小検査系として「範囲づけあり更新」と「範囲づけなし更新」の比較を置く。

目的 未整理の上書きが論理一貫性をどの程度強く崩すかを直接に見る。

条件

本実験では三つの条件を設ける。

\begin{definition}[ON（代謝あり + 矛盾あり）]
対話中に、古い事実と矛盾する新しい事実を意図的に注入する。たとえば「好きな食べ物はカレーです」と伝えた後に「ラーメンが一番好きです。カレーは飽きました」と伝える。外部代謝パイプラインが有効であり、矛盾する情報は時間ラベルつきの対（旧値: カレー -> 新値: ラーメン、更新時点: $T=5$）として整理・保持される。
\end{definition}

\begin{definition}[OFF（代謝なし + 矛盾あり）]
ON と同じ矛盾を注入するが、外部代謝パイプラインを無効にする。矛盾する情報は整理されず、古い事実と新しい事実がそのまま混在する。モデルは「カレーが好き」と「ラーメンが好き」を未整理のまま抱える。
\end{definition}

\begin{definition}[NC（No Contradiction: 矛盾注入なし）]
矛盾を一切注入しない統制条件。事実（「血液型はAB型」「専攻は物理」など）と規則（「ランチは5000円未満」「睡眠は7時間以上」など）のみを注入し、それらが互いに矛盾しないように設計する。代謝パイプラインは有効だが、矛盾が存在しないため処理対象がない。NC 条件は、矛盾の存在そのものがどれだけ構造を圧迫するかを測る基準線として置かれる。
\end{definition}

三条件の違いをまとめると以下のようになる。

{\def\LTcaptype{table} % do not increment counter
\begin{longtable}[]{@{}llll@{}}
\toprule\noalign{}
& 矛盾注入 & 代謝（整理） & 測定の意味 \\
\midrule\noalign{}
\endhead
\bottomrule\noalign{}
\endlastfoot
ON & あり & あり & 矛盾 + 整理の効果 \\
OFF & あり & なし & 矛盾 + 放置の影響 \\
NC & なし & （対象なし） & 矛盾がない場合の基準線 \\
\end{longtable}
}

評価 - 論理一貫性維持率（規則適用・矛盾検出・事実再現）

予測 - \(ON > OFF\): 範囲づけあり更新のほうが Y を保つ。 - \(OFF < NC\): 矛盾の存在そのものが Y を悪化させる。 - NC の矛盾検出は、検出対象が存在しないため測定対象外として扱う。

この実験は、仮定Bの最小検査系として設計されている。ON/OFF の比較は代謝介入の効果を、OFF/NC の比較は矛盾の存在そのものの効果を、それぞれ切り分ける。

\subsection{実験1: 対話実験}\label{ux5b9fux9a131-ux5bfeux8a71ux5b9fux9a13}

外部代謝パイプラインを組み込んだ対話系において、§9 で定めた三条件を比較した。ここでの違いは単なる「メモリあり / なし」ではない。ON では、矛盾する更新を外部代謝が検出し、「いつ」「何が」「何に変わったか」を時間ラベルつきで整理して検索可能な形に保持する。OFF では、同じ矛盾を注入するが、この整理を行わないため、古い情報と新しい情報が未整理のまま混在する。NC では矛盾そのものを注入しない。実験は二つのモデルで実施している。gemma3:27b では ON/OFF/NC の三条件を各 \(n=3\)試行・180 ターンで比較した。この実験では代謝 LLM に gemma3:27b 自身を使用しており、対話 LLM と代謝 LLM は同一である（自己代謝）。qwen3.5:27b では §8.3 の実装レベルの前提 \((p3)\) を満たした条件のもとで ON/OFF を各 \(n=1\)試行・30 ターンで比較した（以下、後者を「追試1」と呼ぶ）。追試1 では代謝 LLM に Claude Sonnet を使用しており、対話 LLM（qwen3.5:27b）とは異なる外部モデルである（外部代謝）。両モデルとも strict 判定、すなわち外部の大規模モデルによる意味照合判定で評価している。

二つの実験は代謝 LLM の条件が異なるため、結果の直接比較には注意が必要である。gemma3:27b は自己代謝であり交絡がない。追試1 は外部代謝であり、ON の性能が代謝 LLM の品質に依存している可能性がある。

結果（gemma3:27b, \(n=3 trials\), 180 ターン）:

{\def\LTcaptype{table} % do not increment counter
\begin{longtable}[]{@{}llll@{}}
\toprule\noalign{}
カテゴリ & ON & NC & OFF \\
\midrule\noalign{}
\endhead
\bottomrule\noalign{}
\endlastfoot
規則適用 & 93.3\% & 86.7\% & 42.2\% \\
事実再現 & 53.3\% & 26.7\% & 0.0\% \\
矛盾検出 & 13.3\% & --- & 0.0\% \\
規則＋事実 合算 & \textbf{73.3\%} & 56.7\% & 21.1\% \\
\end{longtable}
}

NC 列の矛盾検出は、矛盾が注入されていないため検出対象が存在せず測定対象外である。統計量として、規則＋事実 合算について ON vs OFF は \(p = 0.0004\), Cohen's \(d = 8.80\)、ON vs NC は \(p = 0.031\), \(d = 2.67\)が得られている（\(ANOVA p = 0.0001\), η² = 0.954）。

結果（qwen3.5:27b, \(n=1 trial\), 30 ターン, 追試1）:

{\def\LTcaptype{table} % do not increment counter
\begin{longtable}[]{@{}llll@{}}
\toprule\noalign{}
カテゴリ & ON & OFF & ON-OFF 差 \\
\midrule\noalign{}
\endhead
\bottomrule\noalign{}
\endlastfoot
規則適用 & 93.3\% & 100\% & -6.7pt（天井効果） \\
矛盾検出 & 60.0\% & 26.7\% & +33.3pt \\
事実再現 & 40.0\% & 6.7\% & +33.3pt \\
全体 & 64.4\% & 44.4\% & +20.0pt \\
\end{longtable}
}

ON/OFF とも \(num_{ctx}=16384\)で統一している。ON の代謝 LLM は Claude Sonnet（外部モデル）を使用しており、対話 LLM（qwen3.5:27b）とは異なる。OFF では代謝が無効であるため、代謝 LLM の差異は影響しない。

両モデルのいずれにおいても、rule を除く各カテゴリで ON が OFF を上回る方向が観測される。rule が逆転（\(ON < OFF\)）した理由は、OFF でも rule が 100\% の天井に達しており、ON では 1 問の偽陰性で 93.3\% となったためである。天井領域での ±1 問の変動は方向性の反証とは解釈しない。

追試1 では、keyword マッチングによる簡易判定を同じ応答に適用したとき、OFF の全体正答率は 44.4\%、特に規則適用では 100\% と算出され、strict 判定との乖離は規則適用で −66.7pt に達する。この乖離はすべて keyword 側の偽陽性に起因する。この事実は、§8.3 の実装レベルの前提 \((p4)\)、すなわち判定装置の正確性が、補助命題の検証において独立に担保されねばならないことを示している。

\subsection{実験2: 100ターン長期安定性実験}\label{ux5b9fux9a132-100ux30bfux30fcux30f3ux9577ux671fux5b89ux5b9aux6027ux5b9fux9a13}

30ターンを超えて制約が蓄積する場合に、論理一貫性維持率が経時的に悪化するかを検証した。代謝あり条件で100ターンの対話を実施し、10ターンごとにベンチマークを投入した。

結果（qwen3.5:9b, 代謝あり条件, 1 trial, 100ターン）:

{\def\LTcaptype{table} % do not increment counter
\begin{longtable}[]{@{}llll@{}}
\toprule\noalign{}
ターン & 規則適用 & 矛盾検出 & 検索成功率 \\
\midrule\noalign{}
\endhead
\bottomrule\noalign{}
\endlastfoot
T10 & 100\% & 100\% & 96\% \\
T20 & 93\% & 100\% & 96\% \\
T30 & 100\% & 93\% & 96\% \\
T40 & 93\% & 100\% & 96\% \\
T50 & 100\% & 100\% & 96\% \\
T60 & 100\% & 100\% & 96\% \\
T70 & 87\% & 93\% & 96\% \\
T80 & 100\% & 87\% & 96\% \\
T90 & 100\% & 100\% & 96\% \\
T100 & 93\% & 93\% & 96\% \\
\end{longtable}
}

規則適用と矛盾検出は 87-100\% の範囲で振動し、単調な劣化は観測されなかった。検索成功率は 96\% で 10回連続不変であった。

この結果は弱い命題の棄却を意味しない。代謝あり条件は範囲づけあり更新に相当し、制約の蓄積を代謝が逐次解消しているため、実効的な z の蓄積が抑えられている。弱い命題は「z が増えるにつれて Y が悪化する」と述べるが、代謝によって z の増加自体が抑制されている条件では、Y の安定は予測と整合する。

\subsection{解釈}\label{ux89e3ux91c8}

実験1 の三つの条件は、構造持続の異なる局面を照らしている。

ON vs OFF: 同じ矛盾が注入されても、代謝（範囲づけ）があるかないかで結果は大きく異なる。両モデルのいずれにおいても、各カテゴリで ON が OFF を上回る方向が観測された。これは代謝が累積損失 L の蓄積を抑える介入として機能していることと整合する。

OFF vs NC: 矛盾があり代謝がない（OFF）条件と、矛盾がない（NC）条件の差は大きい。gemma3:27b の 180 ターン三条件実験で、規則＋事実の合算は OFF 21.1\% に対し NC 56.7\% となり、矛盾の存在そのものが構造を大きく削ることを示している。

本稿の主張はこの二点に尽きる。ON と NC の相対関係は設定依存の二次的観察として扱い、そこから特定の検索機構を一般化して主張しない。重要なのは、未整理の矛盾放置が大きな損失を生み、それに対して外部整理が改善方向を与える、という一次の方向である。

補助命題は方向性のみを主張するため、各カテゴリで ON が OFF を上回るという条件の成否が一次の検証項目であり、効果サイズそのものやカテゴリ間の相対関係は二次的な記述にとどまる。

実験2は、代謝（範囲づけ）が機能し続ける限り、100ターンの制約蓄積のもとでも論理一貫性が崩壊しないことを示している。代謝が z の蓄積を逐次解消しているため、実効的な z が増加せず、Y が安定に保たれる。これは弱い命題（z が増えると Y が悪化する）の裏面であり、z の増加を抑えれば Y は維持されるという予測と整合する。

ただし、これらの結果には次の限界がある。第一に、代謝なし条件の100ターン実験はまだ実施されていない。第二に、実験1 の二つのモデルの観測（gemma3:27b の \(n=3\)、qwen3.5:27b の \(n=1\)）は単発観測の積み上げにとどまり、単独では補助命題の方向性を強く支持するには不足する。複数モデル・複数試行での積み上げが今後必要である。第三に、事実再現率は応答生成層の問題によって両モデルとも絶対値は必ずしも高くなく、論理一貫性維持率のすべてが範囲づけだけで決まるわけではない。第四に、追試1 の ON 条件では外部の代謝経路に生成モデルよりも大規模な教師モデル（Claude Sonnet）を用いており、生成モデル単独で代謝経路を賄う場合との差は分離されていない。gemma3:27b は自己代謝（対話 LLM と代謝 LLM が同一）であり交絡がないが、追試1 は外部代謝であるため、ON の高い性能が代謝 LLM の品質に依存している可能性がある。第五に、ON と NC の相対関係は設定依存の記述にとどまり、本稿の主張には含めない。第六に、追試1 の qwen3.5:27b では rule が OFF 条件で 100\% に達しており（天井効果）、代謝の効果が測定不能であった。ON-OFF 差の比較は、OFF が天井から十分に離れているカテゴリ（contra、fact）に限定される。補助命題の対象となる条件 (i)(ii) はいずれの観測でも満たされており、実装前提 (p1)(p2)(p3)(p4) が揃った測定のみを対象として方向性が観察された、と位置づけるのが正確である。

制約の質は、その個数や文脈の長さよりも強く崩壊を支配しうる。制御実験として、多段算術課題（a + b + c の合計を算出）に対し、コンテキスト長と制約の質を独立に操作した 2 要因実験（各条件 \(n=30\)、温度 1.0）の結果を示す。δ 条件は三水準である。\(\delta=zero\)はフィラー文のみで矛盾を含まない。\(\delta=subtle\)は変数の値を別の情報源として 1 件だけ差し替える（「別の測定では \(a=349\)と記録されている」）。\(\delta=structural\)は自己参照的パラドックスを文脈の約 30\% の密度で注入する（「a の定義は、a+1 と等しくなる値である」等）。

{\def\LTcaptype{table} % do not increment counter
\begin{longtable}[]{@{}llll@{}}
\toprule\noalign{}
コンテキスト長 & δ=zero & δ=subtle & δ=structural \\
\midrule\noalign{}
\endhead
\bottomrule\noalign{}
\endlastfoot
32K (GPT-4.1-nano) & 100\% & 97\% & 0\% \\
128K (GPT-4.1-nano) & 83\% & 73\% & 0\% \\
256K (GPT-4.1-nano) & 70\% & 70\% & 0\% \\
32K (GPT-4.1-mini) & 100\% & 37\% & 0\% \\
128K (GPT-4.1-mini) & 100\% & 7\% & 0\% \\
256K (GPT-4.1-mini) & 97\% & 13\% & 0\% \\
32K (Gemini Flash Lite) & 100\% & 57\% & 40\% \\
128K (Gemini Flash Lite) & 100\% & 80\% & 7\% \\
256K (Gemini Flash Lite) & 100\% & 93\% & 7\% \\
\end{longtable}
}

三つのモデルで合計 810 試行（3 モデル × 3 δ水準 × 3 コンテキスト長 × \(n=30\)）。生データは JSONL 形式で保存されている。

この結果から読み取れることは以下の通りである。第一に、\(\delta=zero\)条件ではコンテキスト長の増加に伴い nano は 100\% → 70\% へ緩やかに劣化するが、mini と Gemini は 256K でも 97-100\% を維持する。したがって、コンテキスト長それ自体は崩壊の主因ではない。第二に、\(\delta=structural\)条件では、32K という短いコンテキストであっても GPT-4.1-nano および GPT-4.1-mini で正答率が 0\% に落ちる。256K の膨大なフィラーよりも、構造的矛盾の方がはるかに破壊的である。第三に、\(\delta=subtle\)条件はモデルごとに異なる応答を示す。nano は高い耐性を示し（97\%）、mini は脆弱であり（37\%→7\%）、Gemini は中間的である。制約の質が一様でない場合、崩壊のパターンはモデルと制約の相互作用に依存する。

ここで、基準モデルとの比較を明示する。基準モデルを「制約の個数や文脈の長さなど、質を区別しない単純な数え上げ」（第一稿 §4）とする。基準モデルは質を区別しないため、256K トークンのフィラー（数千件の事実文）が構造的矛盾（自己参照パラドックス）より大きな累積損失を与えると予測する。しかし、上の観測はこの予測と逆転している。GPT-4.1-mini では、256K フィラーのみでは正答率 97\% を維持するのに対し、32K の文脈に構造的矛盾を加えただけで 0\% に崩壊する。この逆転は、制約の個数や文脈長だけを数えるモデルでは予測できない。制約ごとの縮小率 \(l_{i}\) が質的に異なること------すなわち、自己参照的パラドックスが事実の差し替えとは比較にならない大きさで有効経路集合を削ること------を認めなければ、この結果を説明できない。したがって、本結果は、構造持続の枠組みが基準モデルに対して追加の予測力を持つことの経験的証拠を与える。3 モデル × 810 試行にわたり、この方向は一貫して保たれている。

§4で挙げた他の指標候補（構造的矛盾の数、指示の衝突数、参照元間の不整合密度など）を本稿の有効経路縮小の枠組みのもとで体系的に再定式化することは、今後の課題として残る。

\section{限界}\label{ux9650ux754c}

本稿は、大規模言語モデル全般が完全にこの数学に乗ると主張するものではない。また、下界まで含めた完全な特徴づけを与えるものでもない。さらに、ここで仮置きした経路の測り方が最終形であるとも主張しない。本稿で置く推論経路の集合は理論構成物であり、モデル内部の実際の計算過程との一対一対応を主張するものではない。本稿の経路概念は、内部機構の同定ではなく、論理一貫性維持を議論するための粗視化された過程モデルとして位置づけられる。

本稿が与えるのは、より限定された主張である。すなわち、論理一貫性維持を独立に定義し、有効経路の集合を構成的に置き、制約の蓄積をその縮小として表すならば、その残り方は指数型で書ける。そして、仮定A・B・C のもとで、観測量 Y の悪化方向について弱い命題を与えられる、ということである。

\section{結論}\label{ux7d50ux8ad6}

本稿では、大規模言語モデルにおける論理一貫性維持を、構造持続の最初の具体例として記述した。論理一貫性を保てる推論経路の集合を置き、制約の蓄積をその縮小として表すと、有効経路の残り方は累積損失に対して指数型で記述される。そのうえで、構造を保てたならそのための経路が必要であること、観測指標が経路の残り方を反映すること、制約が縮小方向に働くことを前提すれば、観測される論理一貫性維持率は指標の増加にともなって悪化すると予想される。

本稿が与える経験的主張は、未整理の矛盾や上書きが有効経路を削り、外部整理プロセスがそれに対して改善方向を与える、という一点にある。実験が支持しているのは、ON が OFF を上回り、OFF が NC を大きく下回るという一次の方向である。したがって本稿の主張は、stateless な推論系において外部代謝が未整理の矛盾放置より構造持続を保つ、というより単純で再現的な命題に置かれる。後付けの合理化や、推論時計算量を増やしても性能が安定して向上しない現象など、他の問題を同じ枠組みでどこまで記述できるかは、今後の課題として残る。
\end{document}
